\label{chapter:nq}

\section{What Question Datasets are There?}

In the chapter on the Cranfield Paradigm (Chapter~\ref{}), we argued
that reusable test collections were the life blood of research: it
lets you know when you've made an advance, facilitating the process of
standing on the shoulders of giants.
%
This chapter will outline the question answering datasets that enabled
the rapid progress of question answering algorithms.

However, I don't want to just uncritically give a listing of all the
datasets that have come out.
%
There are \emph{a lot} of them.

In addition to not being a complete listing, I will also be critical
of datasets, mostly through the lens we developed in the chapter on
the Manchester Paradigm (Chapter~\ref{}).
%
Again, I ask my friends in the research community to not take this too
harshly.
%
And some of these flaws took a while to figure out, and I'm writing
this with the advantage of hindsight!
%
So the selection of datasets will also reflect what I think makes for
the best narrative about what we know about how to make a good
dataset.

Because I'm ordering the story chronologically, the broad trend that
we'll observe is that the amazing tour-de-force of \textit{Watson}
created a renewed interest in Question Answering, there was an
explosion of datasets trying many things, and the resulting consensus
was to adopt the Manchester paradigm once question answering was
extensively used to benchmark \abr{ai}.

But we also need to discuss what \emph{is} a question answering
dataset for the purposes of this chapter.
%
There are many question answering datasets out there that are
essentially logic problems or \abr{sql} query lookups on a limited
domain and we're not going to talk about them.
% todo(jbg): Add some examples
%
That's a very valid line of research and it essentially corresponds to
a form of parsing where you have some text and you need to turn that
into a logical form or an \abr{sql} statement.  If you parse it
correctly, you get the answer.\footnote{You will see questions like
this in other datasets we discuss; this is not a contradiction!
Rather, if your dataset is only about asking Geography questions from a
list of \abr{us} states, that dataset would not be included here.
However, if you could have a question about the second largest city in
Louisiana alongside computing the angles of a triangle, that's fine.}

Instead, inspired by the breadth of topics admitted by the Manchester
paradigm exemplar the Turing Test, we want to consider any question
that can be expressed in natural language.
%
This contains all of the ambiguity, complications, and complexity we
covered in the Manchester chapter.
%
In addition to the language component, it also requires a lot of
knowledge.
%
In these question answering datasets, you need to make sense of a
large body of text to answer questions.  Moreover, you usually need to
be able to answer questions about just about anything, not just a
narrow set of topics.

\section{The first \abr{qa} Datasets}

Instead, we're going to talk about more open-ended question answering
that’s often called ``machine reading'' question answering for reasons
that will be clear in a second.
%
Despite the organization of the book, these datasets developed
alongside algorithms for question answering.
%
The shape of these datasets heavily constrained what algorithms could
be successful at answering them.
%
We'll describe how these algorithms work more in the next chapter, but
the rough outline is that they take a question in as input, find a
document that may answer the question, and then highlight the exact
phrase that answers the question within that document.

Okay, back to the datasets.  Rather than just listing them off, I’m
going to try to organize the datasets based on: how much information
you have to answer the question and how complicated the answer is.
Let’s say that we can write this as a two dimensional plot with
provided information as the $x$ axis and answer complexity as the $y$
axis.  I’ll roughly go chronologically through time.

And then we will put in outline whether the dataset marks a Manchester
or Cranfield dataset.  Several of the early datasets are neither.
Rather, they are created by novices, so I will mark each dataset's
position with an icon: a question mark for information-seeking
Cranfield questions, a buzzer for Manchester questions, and a
pawn\footnote{I use a pawn, the simplest chess piece, to mark the
relative lack of sophistication of the hard working people who put the
dataset together and as a nod to Mechanical Turk.  The original
Mechanical Turk was a sham machine that could purportedly play chess
but actually had a human of diminutive stature inside.  Amazon created
a platform of that name to provide both the technical infrastructure
and the marketplace to find people who could complete tasks that
required human intelligence. } for the novices.
% TODO(jbg): finalize novice symbol.  It could be dunce cap triangle,
% pawn for Mechanical Turk, or something else.

% TODO(jbg): Add this figure if it still makes sense

\section{TREC QA}

Let’s start with the most simple example of a question answering dataset: you get a question as input and you need to produce a short answer.  
%
This is often called factoid \abr{qa}.\footnote{Karl Moritz Hermann really dislikes this term, as in the \abr{uk}, a ``factoid'' is not necessarily true.}
%
The prototypical example of this is \abr{trec qa}; you'll remember \abr{trec} as the Text REtrieval Conference, a conference organized by the \abr{us} National Institute of Standards and Technology (\abr{nist}), the home of Ellen Voorhees, who coined the term ``Cranfield Paradigm''.  
%
This is a super old school dataset, but important to remember, as it set the form of everything that came after.

The actual dataset is tiny: only a couple hundred questions.  
%
And the official evaluation process seems impossible by today’s standards: you had professional annotators judge whether a system provided the right answer.
%
While every other dataset we’ll talk about has more automatic evaluations, this is probably the best way to do it, at least for some trickier answers.

This is our first QA dataset.  The questions are not very difficult, but the answers are more complicated than many of the datasets we'll discuss later.
%
It's a little tricky to put this on the chart because it had multiple iterations.  In some iterations, the organizers wrote the questions.  Which could make it a Manchester dataset, but it might have been closer to crowdworkers, as there wasn't much theory available about how to write a good question.
%
However, in later years the questions were gathered from search engine logs, which would make it a Cranfield dataset.
%
Thus, I'll mark it as a Cranfield dataset.

However, while it was similar to a ``machine reading'' dataset, it was not a true ``machine reading'' dataset like the ones we'll discuss next.
%
While the questions had answers that could be found in the documents, the answers were not trivial to find: the algorithms had to do more than just highlight a substring to produce the answer.
%
Thus, as we move forward in time, we'll see that the simpler datasets enabled building models that could answer more difficult questions once they algorithms ``got their reps in'' on the simpler datasets.

\section{CNN / Daily Mail}

One of the first machine reading data sets was the \abr{cnn} / Daily Mail data set from Karl Moritz Hermann and company.  Here you have a bunch of news articles from online news sources (hence the name).
%
This dataset started the trend of conditioning a question on a single document: in this case a news article.  
%
They asked you to answer questions about those news articles.  For example, \question{who reached the Cricket World Cup}.  Somewhere in the document was the answer \answer{New Zealand} but you need to find that.

One thing that is controversial is that you didn’t actually find the string \answer{New Zealand} but rather an anonymized symbol.
%
This meant that you couldn’t just have a language model memorize all the answers (which is even more of a problem now that it was back then).  However, this had another drawback!  
%
Because there weren’t that many symbols, the task was essentially a multiple-choice question: you identified the entities in the paragraph and then you picked which one was correct.  
%
Despite this problem and and that introduced a little bit of noise to the problem, you can see the influence of this question answering dataset on the rest of the field, particularly on ``the next big thing'': \squad{}.

Here the kind of dataset is clear: it's neither a Manchester dataset nor a Cranfield dataset: the questions were automatically generated from human summaries of the news articles.
%
Thus, the questions are neither what a human would ask nor what would challenge a computer.

While both of these datasets were significant, the next dataset made the biggest splash and arguably supercharged the ``machine reading'' setting we'll discuss in detail in the next chapter.

\section{SQuAD}

The Stanford question answering dataset (\squad{} for short) became the gold standard.  
%
While it might not be the best dataset, it was the most popular for the longest; while it was supplanted by Natural Questions (we'll cover it in depth shortly), \squad{}'s reign was longer.
%
If you didn't report results on \squad{}, you needed to explain why.
%
Not only did this dataset popularize ``machine reading'', it arguable made \abr{qa} an academic subfield, and that subfield established the evaluations that made modern \llm{} evaluation (and thus the underlying models chasing those metrics) possible.  
%
Instead of a book about the connections between question answering and the foundations of intelligence and society, this might have been a phamplet of me complaining about Watson.

However, despite this foundation, the questions are again neither Cranfield nor Manchester. 
%
These questions use paragraphs from Wikipedia ask crowd workers to write questions about that paragraph.  
%
The questions are intended to have an answer that is a span within that Wikipedia paragraph.  
%
One way of thinking about this is that a span is a contiguous span of characters: you highlight that span to give the answer.

This is a very large data set so it doesn't have many of the problems of the smaller data sets that we've talked about before. 
%
The downside of this is that these crowd workers can often be reverse engineered there are tricks and cheats that they use to create questions.  
%
For example, if a question asks ``when'', then all a QA system really has to do is find the thing that is being asked about and then find the closest thing that looks like a date nearby.  
%
In other words, the machine learning algorithms are essentially reverse engineering those annotators rather than solving the intrinsic problem.

% TODO(jbg): Add a SQuAD example with a date

The real reason, in my opinion, that \squad{} has become the gold standard is that it has a very well trafficked leaderboard where people can submit their systems and to see how they stack up against the rest of the world.
%
We'll talk more about how important leaderboards are to the development of \abr{ai} generally, but \squad{} started the trend of reusable test collections available online for \emph{anyone} to use as part of the process of developing new methods.
%
While \abr{trec} created datasets that you could submit to an annual event, Percy Liang's Codalab was unique in that you could upload your system, it would run on \squad{}, and your system would appear on the leaderboard.
%
This lead to a popularization of a scientific dataset.
%
While this is mostly a good thing, it led to some headlines that helped confuse consumers of scientific news.


%
Because the answers are spans within the original document this reduces the task to multiple-choice much like the Daily Mail/CNN data set that we talked before.  
%
Thus, it doesn't allow answers that are latent in the text that aren’t a fixed span.  
%
For example, if the question is \question{What states share the Delmarva peninsula with Maryland} and the Wikipedia page has 
\evidence{The Delmarva peninsula is made up of Deleware, Maryland, and Virginia}, it can’t answer \answer{Delaware and Virigina} because there’s a pesky Maryland in the middle!



So putting this on the chart, I’m going to put SQuAD a little bit above CNN/DailyMail.
%
They’re both implicitly multiple choice, but there are more spans than entities, so answering SQuAD questions is a little more difficult.


\section{Search Engines: The Source of Natural Questions}
\label{sec:datasets:nq}

One criticism of SQuAD questions is that their questions are unnatural.  They’re written by crowdworkers paid to write artificial questions.  Two datasets that use “real” questions are \abr{ms} Marco and Natural Questions. 
%
They’re both somewhat similar, so I’ll mostly talk about Natural Questions, which is newer, as a stand-in for both.  
%
These questions are harder than \squad{} and Trec questions!  That’s because the questions come from people using search engines.
%
These people don’t always know what they’re asking about (that’s why their using a search engine after all), so sometimes they have false assumptions, ambiguities, and other problems (Chapter~\ref{}).  

Another facet of these datasets that make the questions harder is that
search engines---built on the techniques we discussed in
Chapter~\ref{ch:ir}---trained people over years to answer questions
themselves by finding webpages.  Nobody gets good answers when they
type in a question like ``what is the last capital in Western Europe to
be captured by a Muslim force''\footnote{The answer is XXX, but to
figure that out, you are essentially doing a multi-hop query (which we
discuss in Chapter~\ref{ch:formats}), so a human trying to answer this
question would decompose it into multiple separate queries.  
%
Now,
however, with \abr{ai}, you often \emph{will} get answers to these
sorts of questions, although you shouldn't always trust them. } so if
you type that into Google you likely will not get an answer to that
question directly.  Instead, you’ll piece together multiple queries to
answer your question.

In other words, the failure of search engines to answer questions over the last couple of decades limited the usefulness of the kinds of questions that Google could get from its users in 2018.  

So how does Natural Questions find the answers to these questions?
%
They find Wikipedia given the question as a search string: an information retrieval problem!  
%
They then have annotators validate that this is a good search result and then have the annotators highlight the answer string.  Now this isn’t perfect.  

Here are some examples where I disagree with some of the implicit assumptions that get made in the NQ dataset.  
%
For example, it assumes that when you’re asking about ``Michigan'', you’re asking about the ``university of michigan football team'', when you’re asking about ``hockey'', you’re asking about ``men’s ice hockey'', and when you’re asking about ``supreme court'' you’re asking about the ``indian supreme court''.  
%
This is a byproduct of now the \abr{nq} dataset was created: the queries are too short (the users are trained to type only keywords), and the \abr{ir} system did an arbitrary disambiguation.
%
We'll talk more about how this might be problematic as we consider \abr{nq}'s long-term legacy.

% TODO(jbg): Add the figure about going back and forth

% As we put these datasets up on the chart, I want to make a distinction between the ``Traditional'' and ``Open'' interpretations of natural questions.  
%
% While initial natural questions was treated as a SQuAD-like dataset, recent interpretations have not assumed that you know which Wikipedia page had the answer: you also need to do the \abr{ir} component as well.  
%

\section{Early Manchester Datasets}

As you might guess, from the Manchester chapter, I think that trivia writers would do a better job of writing questions to train question answering systems than either crowdworkers or users (\textit{a la} the Cranfield paradigm).
%
Let’s talk about three of them: SearchQA, TriviaQA, and \abr{qanta}.
%
SearchQA and TriviaQA actually look a lot like Natural Questions, but instead of the queries coming from random people on the internet, they come from professional writers.  
%
Researchers go out and find pages given the query, highlight the answers.  
%
I’m not going to dwell too much on these datasets because the copyright is a little murky.

Instead, let's now turn to my hobby horse: quiz bowl, which I introduced in Chapter~\ref{}.
%
As mentioned before, it's played by thousands of trivia nerds every weekend, particularly at the high school and college level.  
%
Unlike SearchQA and TriviaQA, there’s an explicit community consensus that old quiz bowl questions are public domain.  
%
So as a result we get a lot of free data from experts. These are not crowd workers: these are people who are passionate about human question answering, and we're using their expertise and their experience to improve computer question answering.

This is the gold standard in human question answering (i.e., trivia tournaments), and my goal is to make it more popular than it currently is for computer question answering.  
%
I first made the case that we should be using these questions in 2012, and we called the resulting dataset \abr{qanta}, for Question Answering is Not a Trivial Activity.\footnote{In our first paper, it was originally called Question Answering using Neural 
% TODO(jbg): Look up original acronym from paper with mohit
}  
%
I've already made the case why I think this is a good idea in this book, so I won't repeat it here, but I want to talk about the historical context a little bit more.

At the time, I hadn't thought about the Manchester and Cranfield distinction, so I only described what the dataset was like.
%
I thought that its charms would be enough to get people to work on it.  
%
That failed.  Other datasets were obviously much more popular.  
%
Part of it was because our rinkydink Maryland team didn't have the ad budget of Google.  Part of it was that the format was too complex.  And a big part of it was that I didn't make my case well.

I've now spent years thinking about this, and I built up this Manchester / Cranfield distinction because of it.
%
And I think this gives us a way to look at the differences between the datasets discussed in this chapter.
%
We can't use Manchester or Cranfield to the exclusion of the other, and this also helps explain why crowdsourced datasets are increasingly falling out of favor.

\subsection{The Problem with Crowdworkers}
\label{sec:nq:iid}
% Todo (jbg): Get the offending headline

Creating new machine learning systems requires data.  These data need to come
from somewhere, and \abr{nlp} systems are no different!  If you were to ask a
statistician, they would tell you get a sample from an \iid{} distribution:
independent and identically distributed.  Independent means that example
$x_{n}$ does not depend on on any of the previous examples 1 to n-1, indeed,
all of the examples come from the \emph{same} underlying distribution.

But that’s the problem.  Language isn’t like neutrinos spewed from the sun or
water from a geyser.  Language is inherently an interactive process: and what
I’m going to argue is that this has always been a problem for our NLP
datasets, that it’s getting worse with the advent of muppet Models, and that
we need to address the problem head on going forward.

In Chapter~\ref{ch:manchester}, I introduced the idea of \emph{adversarial} datasets.  
%
One of the biggest complaints about adversarial datasets is that they are not \textit{iid}.
%
Sam Bowman outlined this in the ``Case Against Underclaiming''~\cite{}: adversarial data give too negative a picture of what \abr{ai} can do, since the examples are chosen based on models' current capabilities.
%
As we discuss in Chapter~\ref{ch:leaderboards}, I think this is fair, because an absolute number is meaningless, relative numbers make more sense.  And specific datasets highlight skills in targeted domains.
%
We also saw this in the case of the Turing Test: Sterrett argues that we should compare computer deception to human deception.

One of the biggest sources of \abr{nlp} data are crowdworkers, Sam Bowman's first claim to fame was creating the Stanford Natural Language Inference dataset using crowdworkers.
%
These are not \textit{iid}!
%
Crowdworkers are not representative of the global population~\cite{}: by definition they have access to a computer, they’re typically younger, and they’re focused in particular geographic areas.
%
And even if you restrict to US Crowdworkers (because nobody has virtual private networks or virtual credit cards), this distribution does not match the overall \abr{us} population.

And beyond the underlying distribution issue, there’s a quality issue.  
%
Crowdsourced data often has “shortcuts” that allow systems to “cheat”.  We’ve seen this for:
Entailment: systems can detect whether a claim is true or not just by looking at the premise or the hypothesis (when you should really need both). 
%
Grurangan et al show that crowdworkers use the same tricks to create examples: “at least” is often entailed, irrelevant information is often neutral, and negation is often contradicted.  
%
Poliak et al showed that you can get pretty high accuracy by ignoring the premise and characterized the “tells” that often appeared in the data. 
%
And we had our own paper that even when the premise is required, it doesn’t require deep reasoning~\cite{}.
%
For fact checking: like with entailment, some logical operators like negations are more likely associated with refutation.  We talked about this in our adversarial fact checking paper, Fool Me Twice~\cite{}.
%
And we already talked about how Weissenborn et al.\ showed that you can look at the passage and find the first thing that could possibly fit the answer type: ``When'' finds the year, ``who'' finds a person, ``where'' finds a place, etc.  

\subsection{Users Are the Customers: Why Manchester isn't the Right Answer All the Time}

Speaking of question answering, natural questions claim to be ``natural''\dots these are the data from real users, this actually is \textit{iid} right?
%
And I think the answer is no!  This is because all language is a negotiation—a give and take between two parties—and Google is a part of that negotiation. 
%
Especially in the smartphone era, there's a particular kind of wierdo who uses Android (more associated with Google), not everyone has a smartphone, and not everybody uses Google for answering questions.
%
And even among those using Google, you may not use Google for everything and you might be embarassed to ask Google specific questions.

Apart from who uses Google, Google itself affects the questions it gets.
%
Google is a fantastic tool, but it has its strengths and weaknesses.  
%
It’s really good at looking up who played role X in movie Y.  
%
IMDB (and other structured knowledge bases) make answering this kind of question easy … so people know that if they ask this kind of question from Google they’re likely to get a good answer.  
%
In general, if I can clearly specify entity X and a relationship f(X) that’s going to appear in some KB, then I’m going to get a great answer.

So this means that people are more likely to ask this kind of question.  
%
But longer questions are just harder to answer.  
%
People want to get answers to their questions, so Google has trained users over the last 25 years to ask the kinds of things it can answer (and the average length of queries is going up!).  
%
But at the end of the day, the kinds of questions that Google isn’t so good at are not going to be asked so they’re not going to be a part of the dataset.  
%
And that means that it won’t be part of the all important natural questions leaderboard.  Which means that people will ignore the questions that Google isn’t already good at.

Now, Jordan, aren’t you just being a meanie by picking on Natural Questions?  
%
Your favorite QA dataset that you keep pushing down everybody’s throat—quiz bowl—ain’t iid either.  

And that’s absolutely right!  I like these crazy quiz bowl questions where the distribution has been chosen by hand to reflect the undergraduate curriculum at American Universities. 
%
This also isn’t \textit{iid}, but I think specifying the distribution and being able to defend it is a good thing!

And in addition to the “kinds” of questions in Natural Questions, Google has trained users to ask questions in a particular way: instead of asking \question{who is Amal Alamuddin’s husband}, people type “Amal Alamuddin husband” because Google started as a search engine that ignored function words.  So this also affects what gets into natural questions: Google has trained people to write terse queries, NQ subselects from that, and that becomes the leaderboard.

From a socio-technical perspective, Google is to the turn of the century like muppet models are to today [link in the description if you don’t know why I call things like GPT Muppet Models].  
%
People are learning how to skillfully prompt muppet models, muppet models are impersonating humans on crowdworking platforms (another reason that datasets from crowdworkers aren’t iid), people are using them to help write their essays (which will in turn be used to train the next generation of muppet models when those essays get posted on the web).  

So what does this mean?  
%
Okay, so I’m going to get on my soapbox for a bit … we need to embrace the interaction between users and systems.  
%
And that means something like adversarial examples (Chapter~\ref{ch:manchester}), and it is true that the accuracy numbers are going to be lower than “normal” datasets. 
%
But I think that this makes a category error!  There’s no such thing as a normal, \textit{iid} dataset. 
%
All datasets have their own pathologies, and you cannot readily compare \squad{} numbers to \abr{nq} numbers: there are fundamentally different underlying distributions and data collection practices.
%
The same thing is true for my beloved quiz bowl format.

And I think this is a healthy development.  
%
Rather than obsessing about rank on a leaderboard and the single floating point number with your accuracy that comes out, we should focus on what are the abilities (and weaknesses) of models.  
%
And accepting that there is no such thing as the \textit{iid} bunny bringing you a basket of nicely distributed questions will help us get there faster.

Likewise, it's dangerous to ignore distributions when you compare models or compare models with humans.
%
One thing that really grinds my gears about \squad{} is that the \squad{} dataset provides a ``human upper bound''
(and to be clear this is not the fault of the people who designed \squad{} but from people who misinterpret what they did) 
but there less well informed news organizations or promoters of research that mischaracterize what it means to be near the human bound.  It also mischaracterizes what it means to ``read'' a document.  
%
If machine learning systems are getting close to this human bound it doesn't mean that computers are reading better than humans.  
%
Rather, it means that computers are doing a better job of highlighting particular passages better than underpaid crowd workers.  
%
As you can guess from the Watson chapter (Chapter~\ref{}), this offends my sense of fair play.  I lay out my full theory of what would make for a fun and fair human computer--\abr{ai} game show in Chapter~\ref{}.
%
For this, I think the Manchester paradigm is the much better choice.

That said, the transition from Manchester-style spectacle to Cranfield-style
product is hard. \abr{ibm} Watson's \jeopardy{} victory was a tour de
force of the Manchester paradigm, but turning that success into a
user-facing product was much more difficult~\citep{chu2002piquant,
deutscher-21,jennings2011ibm}. The same tension shows up in dataset
design. \squad{} questions were written by people who knew the
answers, but the data do not follow Manchester-style question-writing
norms; the result was a proliferation of shortcuts and artifacts that
later analyses exposed~\citep{weissenborn-17}. In contrast,
Manchester-oriented benchmarks such as the Winograd challenge
explicitly try to avoid cheap tricks~\citep{levesque2014behavior}. 

The
point is not that one paradigm is superior, but that each requires
different care.
%
Cranfield data are still necessary because the data distribution of uninformed users in the Cranfield paradigm is categorically different from the hyper-informed nerds in the Manchester paradigm.  
%
You ideally need to master both distributions.
%
This is the most enduring lesson of the \abr{nq} dataset years after its release.  Even though it too has fallen by the wayside as newer datasets came to the fore, it still offers important lessons about how real people ask questions, how systems trained on those imperfect questions can go wrong, and what we can do about it.


\section{Natural Questions Five Years Later}
\label{sec:nq:problems}

Natural Questions came out in 2019.  It’s hard to overstate its impact on Question Answering (QA).  Perhaps only second to SQuAD.  
%
\squad{} introduced the extractive format, but \squad{} was “fake” … Natural Questions had true complexity and verisimilitude: people ask crazy questions, some of the questions were unanswerable, and the answers were tied to Wikipedia pages. 
%
And it has yet to be eclipsed … and perhaps it never will.  

Part of the reason that this will probably never happen again is that Natural Questions has two very unique characteristics: questions from real users expensively verified by expensive annotators.  
%
While it’s possible that we might get something like this again in part---\abr{uw} folks crawled Google $n$-grams to extract questions~\cite{}, we paid QuizBowlers to find passages to answer questions~\cite{}---getting them both at once is probably not going to happen again.  
%
But I’d love to be proved wrong! 

However, Natural Questions aren’t perfect.  
%
Sewon Min and company created \camelabr{Ambig}{qa}~\cite{} that showed some of the questions in Natural Questions were ambiguous.
%
Here are some examples that Ben and I talked about in our 2020 paper (and I talked about at more length in the QA datasets lecture, link in the description).  
%
There are multiple interpretations to these questions, but Natural Questions only says that one of them is correct.  

How is the “correct” interpretation chosen in the original Natural Questions paper? 
%
The annotators saw search results from an information retrieval system, and if they found an answer there, they called it the correct answer.  
%
So whatever process biases search results toward English-speaking judges, men’s sports, or flagship universities gets recapitulated in Natural Questions.

Another development is the realization that not all questions are valid and shouldn’t be answered.  
%
There are multiple aspects of this: appropriately calibrating your \abf{qa} system to abstain when it doesn’t know the answer for example.  
%
Natural Questions actually does a pretty good job of supporting this by having unanswerable questions, but it doesn't give a template for \emph{how} to push back against a bad question.

And perhaps we don’t need that any more.
%
As we will see in Chapter~\ref{}, increasingly \abr{ai} models are created from a melange of smaller datasets.
%
Perhaps we don’t need a single monolithic \abr{qa} training dataset to rule them all.  
%
We just need robust evaluation protocols (which can be much smaller). 

And this is, again, a weakness of Natural Questions\dots it’s not a great evaluation set. 
%
I already shared my annoyances with comparisons between human and computer ability on \squad{}.
%
When I was at Google, I tried to create a human--computer competition that was fair in 2020 using \abr{nq}.
%
% I wanted to give the agents---either human or computer---answering the questions comparable information.
%
We called this \camelabr{Efficient}{qa}, and the goal was to see what was the smallest question asnwering system that could work on \abr{nq} questions.
%
Patrick Lewis, the first author of the winning system, took advantage of duplicate information in the dataset (many of the questions were asked again and again) and created a tiny system only dozens of megabytes that clearly won the competition.  It also presaged many of the issues people discovered with \abr{nq}.
%
% To create an equivalence between humans and computers, we let humans answer the questions without seeing any resources 
% Obviously, the open interpretation makes it much harder to find the answer.


Lewis et al. discovered that many QA datasets, including Natural Questions, have a huge overlap between the training set and the testing set. 
%
I’m a Professor, so let’s put this in a classroom context: this is like giving out two thirds of the answers and one third of the questions as a study guide before the test.  
%
Now, it’s a great sanity check that systems can answer those questions, but that means that only a third of the questions are part of a functional test set.  
%
Add that to the problems of ambiguity, and far less of the test set is actually usable, making it difficult to know if we’re making progress on the dataset.

This is something that the trivia community has known for a while.  
%
You can’t reuse clues!  For important test sets (like national trivia tournaments), you need to make sure that you’re finding new, interesting clues that haven’t been used before (or at least for a while).  
%
This is why we’ve always used the flagship national trivia tournaments as our test sets for quiz bowl.

I want to make it clear that \emph{Natural Questions remains a great resource}.  It hasn’t been matched.  I still use it.  I’m jealous of the people who got to build it.  But it’s important to know the limitations. 

Now, for all of the reasons I talked about before, it probably will never be done again.  Nobody is crazy enough to do it. 
%
But how would I do it if someone were foolish enough to give me the resources to create the next iteration of something like it?

Back in the early 2000s, there was a QA service called Aardvark.  
%
The idea was that it would connect up real humans with topics those smart humans could answer questions about.  Question comes in and it gets routed to that person.  
%
That was a great idea but before it’s time.  What you’d do today is build up a large collection of these people (and pay them well), a question from a major search engine comes in\dots we try to detect if the computer can answer it directly.  
%
If yes, we answer it, if not, we farm it out to the experts.  We’ve trained the experts to interrogate the question askers to remove ambiguities, false presuppositions, etc.  
%
To create a “clean” version of the question and then a strong reference answer supported by relevant information.

And to be clear, that’s not what happens if you ask this to ChatGPT\dots it doesn’t really probe the user’s intent but tries to do some correction (not invented by a linguist) and provide the answer to the most likely question.  
%
But still erases all of the people whose work Edison built on.
